\section{مقدمه}

در سال‌های اخیر، توسعه فناوری‌های ارتباطی و گسترش سامانه‌های حمل‌ونقل هوشمند \LTRfootnote{\lr{Intelligent Transportation Systems (ITS)}} زمینه را برای ظهور نسل جدیدی از خودروهای متصل و خودران فراهم کرده است. در این میان، شبکه اقتضایی خودرویی \LTRfootnote{\lr{Vehicular Ad-hoc Network (VANET)}}، به‌عنوان زیرمجموعه‌ای از شبکه‌های اقتضایی سیار \LTRfootnote{\lr{Mobile Ad-hoc Networks (MANET)}}، ارتباط میان خودروهای نزدیک به یکدیگر و ارتباط میان خودروها و زیرساخت‌ها را تسهیل می‌کند \cite{hozouriOverviewVANETVehicular2023}. هدف اصلی از ایجاد چنین شبکه‌ای، فراهم‌کردن بستری برای تبادل سریع و قابل اعتماد اطلاعات است؛ اطلاعاتی که می‌تواند برای افزایش ایمنی، مدیریت ترافیک، ارائه خدمات ایمنی و کمک به راننده و پشتیبانی از تصمیم‌گیری خودروها به کار رود \cite{hozouriOverviewVANETVehicular2023,luConnectedVehiclesSolutions2014}.

اهمیت این فناوری را می‌توان از منظر ایمنی جاده‌ای نیز بررسی کرد. بر اساس داده‌های سازمان بهداشت جهانی، سالانه حدود ۱.۱۹ میلیون نفر در سراسر جهان بر اثر تصادفات جاده‌ای جان خود را از دست می‌دهند و این تصادفات یکی از عوامل اصلی مرگ‌ومیر کودکان و جوانان به شمار می‌رود \cite{farsimadanReviewSecurityChallenges2025,VanetSecurityPrivacy}. از سوی دیگر، بخش بزرگی از تصادفات به خطای انسانی مرتبط است؛ در منبع مروری مورد بررسی، ۹۲ درصد تصادفات به‌طور متداول به دو نوع خطای انسانی نسبت داده شده است: خطاهای قضاوت، مانند بی‌توجهی، نظارت ناکافی و حواس‌پرتی راننده، و خطاهای تصمیم‌گیری، مانند رانندگی با سرعت بیش از حد، واکنش دیر به رویدادها، رعایت نکردن فاصله ایمنی و خطا در برآورد وضعیت \cite{hozouriOverviewVANETVehicular2023}. در چنین شرایطی، توانایی خودروها برای دریافت و ارسال اطلاعات مربوط به وضعیت جاده، موقعیت خودروهای اطراف، شرایط ترافیکی و رویدادهای خطرناک می‌تواند نقش مهمی در کاهش زمان واکنش راننده و افزایش ایمنی داشته باشد \cite{farsimadanReviewSecurityChallenges2025}.

یکی از مهم‌ترین قابلیت‌های شبکه اقتضایی خودرویی، انتشار هشدارهای ایمنی در زمان نزدیک به وقوع رویداد است. برای نمونه، خودرو می‌تواند اطلاعاتی درباره ترمز اضطراری، اجتناب از برخورد، هشدار برخورد در تقاطع، بسته بودن جاده، شرایط آب‌وهوایی و هشدار ازدحام جاده‌ای را به خودروهای مجاور منتقل کند \cite{hozouriOverviewVANETVehicular2023}. چنین پیام‌هایی توسط واحد سوارشونده \LTRfootnote{\lr{On-Board Unit (OBU)}} خودرو دریافت و اعتبارسنجی شده و از طریق سامانه ارتباطات اختصاصی کوتاه‌برد \LTRfootnote{\lr{Dedicated Short-Range Communications (DSRC)}} برای خودروهای دیگر پخش می‌شوند \cite{hozouriOverviewVANETVehicular2023}. این قابلیت موجب می‌شود خودروهای حاضر در یک محدوده تنها به اطلاعات حسگرهای درونی خود متکی نباشند و بتوانند از اطلاعات تولیدشده توسط سایر خودروها و زیرساخت‌های جاده‌ای نیز بهره ببرند \cite{xuInternetVehiclesBig2018}. در نتیجه، شبکه اقتضایی خودرویی می‌تواند به‌عنوان لایه ارتباطی مکملی برای سامانه‌های ایمنی و کمک‌راننده عمل کند و زمینه را برای توسعه کاربردهای پیشرفته‌تر، از جمله رانندگی خودکار، فراهم سازد؛ داده‌های گستردهٔ اینترنت خودروها \LTRfootnote{\lr{Internet of Vehicles (IoV)}} فناوری کلیدی برای خودروهای خودران تلقی می‌شوند \cite{xuInternetVehiclesBig2018}.

ارتباطات در این محیط محدود به ارتباط مستقیم میان دو خودرو نیست. معماری شبکه اقتضایی خودرویی می‌تواند شامل ارتباط خودرو با خودرو \LTRfootnote{\lr{Vehicle-to-Vehicle (V2V)}}، ارتباط خودرو با زیرساخت \LTRfootnote{\lr{Vehicle-to-Infrastructure (V2I)}}، ارتباط زیرساخت با زیرساخت \LTRfootnote{\lr{Infrastructure-to-Infrastructure (I2I)}}، ارتباط خودرو با عابر پیاده \LTRfootnote{\lr{Vehicle-to-Pedestrian (V2P)}}، ارتباط خودرو با موانع کنار جاده \LTRfootnote{\lr{Vehicle-to-Barrier (V2B)}}، ارتباط خودرو با ابر \LTRfootnote{\lr{Vehicle-to-Cloud (V2C)}}، ارتباط خودرو با حسگر \LTRfootnote{\lr{Vehicle-to-Sensor (V2S)}} و ارتباط خودرو با پهپاد \LTRfootnote{\lr{Vehicle-to-UAV (V2U)}} باشد \cite{hozouriOverviewVANETVehicular2023}. برخی منابع به این فهرست ارتباط خودرو با شبکه \LTRfootnote{\lr{Vehicle-to-Network (V2N)}} را نیز افزوده‌اند \cite{farsimadanReviewSecurityChallenges2025}. ترکیب این ارتباطات محیطی ایجاد می‌کند که در آن حجم قابل توجهی از اطلاعات میان خودروها، واحدهای کنارجاده‌ای \LTRfootnote{\lr{Road-Side Unit (RSU)}}، سامانه‌های ابری و سایر اجزای زیرساخت حمل‌ونقل مبادله می‌شود. همین افزایش سطح ارتباط، اگرچه قابلیت‌های جدیدی ایجاد می‌کند، سطح مواجههٔ سامانه با حمله و نیازمندی‌های امنیتی آن را نیز افزایش می‌دهد \cite{hozouriOverviewVANETVehicular2023,farsimadanReviewSecurityChallenges2025}.

در واقع، ماهیت شبکه اقتضایی خودرویی تفاوت قابل توجهی با بسیاری از شبکه‌های سنتی دارد. خودروها گره‌هایی با تحرک بالا هستند و موقعیت آن‌ها در مدت زمان کوتاهی تغییر می‌کند؛ در نتیجه، توپولوژی شبکه به‌صورت مداوم در حال تغییر است و ارتباط میان گره‌ها ممکن است به سرعت برقرار یا قطع شود \cite{luConnectedVehiclesSolutions2014,hozouriOverviewVANETVehicular2023}. علاوه بر این، پیام‌های ایمنی معمولاً باید با تأخیر بسیار کم منتقل و پردازش شوند؛ زیرا خدمات ایمنی باید با تأخیر کم و قابلیت اطمینان بالا در دسترس باشند و هر خطا در ارتباط می‌تواند به تأخیر در رسیدن هشدار و در نهایت به تصادف منجر شود \cite{hozouriOverviewVANETVehicular2023}. از سوی دیگر، تعداد زیاد خودروهای حاضر در یک محیط شهری یا بزرگراهی می‌تواند موجب ایجاد ازدحام در کانال ارتباطی، افزایش تأخیر و کاهش گذردهی شود \cite{hozouriOverviewVANETVehicular2023,xuInternetVehiclesBig2018}.

این ویژگی‌ها سبب می‌شوند تأمین امنیت در شبکه اقتضایی خودرویی صرفاً به استفاده از یک الگوریتم رمزنگاری یا یک سازوکار تصدیق هویت محدود نشود. یک سامانه امنیتی مناسب باید بتواند به‌صورت هم‌زمان محرمانگی، یکپارچگی، تصدیق هویت، دسترس‌پذیری، انکارناپذیری \LTRfootnote{\lr{Non-repudiation}} و حریم خصوصی را تأمین کند \cite{farsimadanReviewSecurityChallenges2025,VanetSecurityPrivacy}. از این میان، برخی الزامات ممکن است در ظاهر با یکدیگر در تضاد باشند. برای نمونه، از یک طرف لازم است فرستنده یک پیام ایمنی معتبر قابل شناسایی باشد تا خودروهای دیگر بتوانند به پیام اعتماد کنند؛ از طرف دیگر، افشای دائمی هویت واقعی خودرو می‌تواند امکان ردیابی حرکات و ایجاد نمایه رفتاری از راننده را فراهم کند. بنابراین، طراحی سازوکارهای امنیتی در این محیط نیازمند ایجاد توازن میان قابلیت اعتماد، پاسخ‌گویی \LTRfootnote{\lr{Accountability}} و حفظ حریم خصوصی است \cite{VanetSecurityPrivacy,rahmawatiagustinaSystematicLiteratureReview2025}.

یکی از مهم‌ترین تهدیدهای این حوزه، امکان انتشار پیام‌های جعلی یا دستکاری‌شده است. مهاجم می‌تواند با جعل هویت یک خودرو یا ایجاد پیام‌های نادرست، اطلاعاتی درباره موقعیت خودرو، وضعیت ترافیک یا وقوع یک حادثه منتشر کند. چنین حملاتی در یک شبکه خودرویی صرفاً یک تهدید اطلاعاتی محسوب نمی‌شوند، بلکه می‌توانند تصمیم‌های واقعی خودروها و رانندگان را تحت تأثیر قرار دهند \cite{farsimadanReviewSecurityChallenges2025}. حملاتی مانند جعل پیام، جعل هویت، حمله مرد میانی \LTRfootnote{\lr{Man-in-the-Middle (MitM)}}، حمله بازپخش \LTRfootnote{\lr{Replay attack}} و دستکاری داده‌های حسگر از جمله تهدیدهایی هستند که می‌توانند صحت و اعتمادپذیری اطلاعات مبادله‌شده را کاهش دهند \cite{farsimadanReviewSecurityChallenges2025,VanetSecurityPrivacy}.

در کنار تهدیدهای مربوط به صحت و اصالت پیام‌ها، حریم خصوصی کاربران نیز یکی از چالش‌های اساسی شبکه اقتضایی خودرویی است. خودروها در جریان ارتباطات خود اطلاعاتی درباره موقعیت جغرافیایی، مسیر حرکت، زمان تردد و الگوی رفتاری تولید یا منتقل می‌کنند. در صورتی که این اطلاعات بتوانند به یک هویت ثابت مرتبط شوند، مهاجم می‌تواند با تحلیل داده‌های ارتباطی، مسیر حرکت یا رفتار یک خودرو را در طول زمان دنبال کند. حملات ردیابی موقعیت، پیوندپذیری \LTRfootnote{\lr{Linkability}} و نمایه‌سازی \LTRfootnote{\lr{Profiling}} از جمله تهدیدهایی هستند که ضرورت استفاده از سازوکارهای اختصاصی حفظ حریم خصوصی را نشان می‌دهند \cite{farsimadanReviewSecurityChallenges2025,rahmawatiagustinaSystematicLiteratureReview2025}.

برای کاهش این خطر، استفاده از شبه‌نام \LTRfootnote{\lr{Pseudonym}} یکی از رویکردهای رایج در طراحی سامانه‌های امنیتی شبکه‌های خودرویی است. در این رویکرد، خودرو به جای استفاده دائمی از هویت واقعی خود، از شناسه‌های موقتی استفاده می‌کند که هویت واقعی را پنهان می‌کنند و مشارکت ناشناس را در شبکه ممکن می‌سازند \cite{rahmawatiagustinaSystematicLiteratureReview2025}. هدف از تغییر دوره‌ای این شناسه‌ها، قطع رابطه میان موقعیت‌های قبلی و بعدی خودرو و دشوار‌کردن پیوند دادن پیام‌های مختلف به یک هویت ثابت است \cite{luPseudonymChangingSocial2012}. با این حال، استفاده از شبه‌نام به تنهایی مسئله حریم خصوصی را به‌طور کامل حل نمی‌کند؛ زیرا ورودی‌های ثابت تولید شبه‌نام و قابل پیش‌بینی بودن بازه‌های تغییر آن، و همچنین استفاده مهاجم از اطلاعات جانبی، الگوهای مکانی و زمانی و ویژگی‌های ترافیک، می‌تواند پیوند دادن شبه‌نام‌های مختلف را آسان کند \cite{rahmawatiagustinaSystematicLiteratureReview2025,luPseudonymChangingSocial2012}. بنابراین، طراحی چرخه کامل تولید، استفاده، تغییر و ابطال شبه‌نام‌ها همچنان یکی از مسائل مهم پژوهشی محسوب می‌شود؛ در مرور نظام‌مند مورد بررسی، تنها حدود ۱۲ درصد از ۱۱۳ مطالعهٔ بررسی‌شده همهٔ مراحل این چرخه را پیاده‌سازی کرده بودند \cite{rahmawatiagustinaSystematicLiteratureReview2025}.

از سوی دیگر، یکی از محدودیت‌های رویکردهای سنتی امنیتی، وابستگی قابل توجه آن‌ها به مراجع مرکزی است. در معماری‌های مبتنی بر زیرساخت کلید عمومی \LTRfootnote{\lr{Public Key Infrastructure (PKI)}}، مرجع قابل اعتماد \LTRfootnote{\lr{Trusted Authority (TA)}} نقش مهمی در مدیریت و امنیت ارتباطات و صدور اعتبارنامه‌هایی مانند گواهی دیجیتال و شبه‌نام دارد \cite{rahmawatiagustinaSystematicLiteratureReview2025}. اگرچه این معماری می‌تواند امکان تصدیق هویت و اعتبارسنجی پیام‌ها را فراهم کند، اما وابستگی به یک مرجع مرکزی می‌تواند مسئله نقطه شکست منفرد \LTRfootnote{\lr{Single Point of Failure (SPoF)}} را ایجاد کند؛ از دسترس خارج شدن این مرجع بر همهٔ موجودیت‌های شبکه اثر می‌گذارد و از آنجا که هویت واقعی همه را می‌داند، هدف جذابی برای نفوذ و نقض گسترده حریم خصوصی می‌شود \cite{rahmawatiagustinaSystematicLiteratureReview2025}. چالش‌های مربوط به مقیاس‌پذیری، مدیریت گواهی و ابطال سریع اعتبارنامه‌ها نیز در همین زمینه مطرح است \cite{farsimadanReviewSecurityChallenges2025,hozouriOverviewVANETVehicular2023}.

علاوه بر این، رمزنگاری و تصدیق هویت به تنهایی نمی‌توانند تمام تهدیدهای شبکه اقتضایی خودرویی را برطرف کنند. سازوکارهای پیشگیرانهٔ رمزنگاری در برابر حملات داخلی گره‌های مجاز آسیب‌پذیرند؛ یک گره موذی ممکن است دارای گواهی معتبر باشد و پیام‌های خود را از نظر رمزنگاری به‌درستی امضا کند، اما محتوای پیام‌های آن عمداً نادرست باشد \cite{farsimadanReviewSecurityChallenges2025}. به همین دلیل، در سال‌های اخیر موضوع مدیریت اعتماد و شهرت به‌عنوان مکمل سازوکارهای رمزنگاری مورد توجه قرار گرفته است. در این رویکرد، رفتار گره‌ها بررسی شده و بر اساس سابقه عملکرد، توصیه سایر گره‌ها یا تحلیل رفتار، یک امتیاز اعتماد یا شهرت برای آن‌ها محاسبه می‌شود \cite{razaBlockchainBasedReputationTrust2024,ghovanlooyghajarSBTMSScalableBlockchain2021}.

فناوری دفترکل توزیع‌شده \LTRfootnote{\lr{Distributed Ledger Technology (DLT)}} و به‌ویژه زنجیره بلوکی \LTRfootnote{\lr{Blockchain}} نیز به‌عنوان یکی از راهکارهای مطرح برای کاهش وابستگی به مراجع مرکزی و ایجاد یک بستر قابل اعتماد مورد توجه قرار گرفته است. ویژگی‌هایی مانند توزیع‌شدگی، تغییرناپذیری و امکان ثبت قابل بررسی رویدادها می‌توانند در مدیریت اعتماد، ثبت سوابق امنیتی، مدیریت هویت و جلوگیری از دستکاری اطلاعات مفید باشند \cite{razaBlockchainBasedReputationTrust2024,fengBPASBlockchainAssistedPrivacyPreserving2020}. با این حال، استفاده مستقیم از زنجیره بلوک‌های سنتی در محیط‌های خودرویی با محدودیت‌هایی مانند تأخیر تأیید تراکنش، مقیاس‌پذیری پایین، هزینه تراکنش و مصرف انرژی مواجه است \cite{razaBlockchainBasedReputationTrust2024,ghovanlooyghajarSBTMSScalableBlockchain2021,IOTATANGLE20}. این محدودیت‌ها اهمیت بررسی ساختارهای جایگزین و سبک‌تر را افزایش داده‌اند.

در این زمینه، تنگل IOTA \LTRfootnote{\lr{IOTA Tangle}} به‌عنوان یک دفترکل توزیع‌شده مبتنی بر گراف جهت‌دار غیرمدور \LTRfootnote{\lr{Directed Acyclic Graph (DAG)}} مورد توجه قرار گرفته است. ساختار تنگل با معماری متفاوت خود نسبت به زنجیره بلوکی سنتی، امکان بررسی راهکارهایی با هدف کاهش هزینه تراکنش و بهبود مقیاس‌پذیری را فراهم می‌کند \cite{IOTATANGLE20,silvanoIotaTangleCryptocurrency2020}. از این رو، بررسی کاربرد تنگل IOTA در مدیریت اعتماد، اعتبارسنجی اطلاعات و مدیریت هویت در شبکه‌های خودرویی می‌تواند یکی از مسیرهای مهم پژوهشی باشد \cite{IOTATANGLE20,silvanoIotaTangleCryptocurrency2020,feraudoDIVADIDbasedReputation2024}.

هم‌زمان با توسعه این فناوری‌ها، مفاهیم جدیدتری مانند هویت خودمختار \LTRfootnote{\lr{Self-Sovereign Identity (SSI)}}، شناسه‌های غیرمتمرکز \LTRfootnote{\lr{Decentralized Identifiers (DID)}} و اعتبارنامه‌های تأییدپذیر \LTRfootnote{\lr{Verifiable Credentials (VC)}} نیز برای حل مسائل مدیریت هویت و حفظ حریم خصوصی مطرح شده‌اند. هدف این رویکردها آن است که کنترل بیشتری بر هویت دیجیتال در اختیار صاحب هویت قرار گیرد و اطلاعات هویتی تنها در حد نیاز و با کنترل مناسب افشا شود \cite{mazzocca2025survey}. ترکیب این مفاهیم با شبه‌نام‌ها و دفترکل‌های توزیع‌شده می‌تواند زمینه ایجاد معماری‌هایی را فراهم کند که در آن تصدیق هویت، پاسخ‌گویی و حفظ حریم خصوصی به‌صورت هم‌زمان مورد توجه قرار گیرند \cite{mazzocca2025survey,feraudoDIVADIDbasedReputation2024}.

از منظر دیگری، رشد حجم داده‌های تولیدشده توسط خودروها، پیچیدگی حملات و پویایی محیط شبکه، استفاده از روش‌های هوش مصنوعی و یادگیری ماشین \LTRfootnote{\lr{Machine Learning (ML)}} را نیز به یکی از موضوعات مهم پژوهشی تبدیل کرده است. روش‌های هوشمند می‌توانند برای تشخیص ناهنجاری، شناسایی رفتارهای غیرعادی، ارزیابی اعتماد گره‌ها و تشخیص برخی حملات مورد استفاده قرار گیرند \cite{farsimadanReviewSecurityChallenges2025}. با این حال، استفاده از هوش مصنوعی نیز بدون چالش نیست و مسائلی مانند کیفیت و در دسترس نبودن داده‌های آموزشی، حملات تخاصمی، مسموم‌سازی داده و حفظ حریم خصوصی باید در طراحی چنین سامانه‌هایی مورد توجه قرار گیرد \cite{farsimadanReviewSecurityChallenges2025}.

بنابراین، امنیت شبکه اقتضایی خودرویی را نمی‌توان به یک لایه یا یک فناوری خاص محدود کرد. یک معماری امنیتی جامع باید مجموعه‌ای از سازوکارهای رمزنگاری، تصدیق هویت، مدیریت کلید، حفظ حریم خصوصی، مدیریت اعتماد، تشخیص ناهنجاری، مدیریت هویت و در صورت نیاز فناوری‌های دفترکل توزیع‌شده را در کنار یکدیگر در نظر بگیرد. مهم‌تر از آن، این سازوکارها باید با محدودیت‌های ذاتی محیط خودرویی، از جمله تحرک بالا، تغییرات سریع توپولوژی، نیاز به تأخیر پایین، حجم بالای پیام‌ها و الزامات مقیاس‌پذیری سازگار باشند \cite{farsimadanReviewSecurityChallenges2025,rahmawatiagustinaSystematicLiteratureReview2025}.

در چنین شرایطی مسئله اصلی این گزارش بررسی معماری‌های مختلفی‌ست که به ابعاد مختلف امنیت و حریم خصوصی در شبکه اقتضایی خودرویی می‌پردازند و طرح‌هایی را بررسی می‌کنند که میان امنیت، حریم خصوصی، اعتماد، کارایی و مقیاس‌پذیری تعادل ایجاد می‌کنند. این مسئله از آن جهت اهمیت دارد که افزایش امنیت نباید موجب ایجاد تأخیر غیرقابل قبول در پیام‌های ایمنی شود و حفظ حریم خصوصی نیز نباید به گونه‌ای باشد که امکان پاسخ‌گویی در برابر رفتارهای موذی از بین برود. به بیان دیگر، هدف یک معماری مناسب آن است که خودروهای معتبر بتوانند بدون افشای غیرضروری هویت واقعی خود با یکدیگر ارتباط برقرار کنند، پیام‌های معتبر را از پیام‌های موذیانه تشخیص دهند و در صورت مشاهده رفتار غیرقانونی یا موذی، امکان شناسایی و برخورد با عامل مربوطه وجود داشته باشد \cite{rahmawatiagustinaSystematicLiteratureReview2025,VanetSecurityPrivacy}.

بر این اساس، این گزارش به بررسی جامع ابعاد مختلف امنیت و حریم خصوصی در شبکه‌ی اقتضایی خودرویی می‌پردازد. ابتدا در فصل ۲ مفهوم، معماری، اجزای اصلی و ویژگی‌های این شبکه معرفی می‌شود و در فصل ۳ چالش‌های شبکه‌سازی، مدیریت منابع، مقیاس‌پذیری و استانداردسازی آن بررسی می‌گردد. فصل ۴ الگوی تهدید و طیف حملات امنیتی علیه این شبکه‌ها را طبقه‌بندی می‌کند. پس از آن، راهکارهای دفاعی به تفکیک بررسی می‌شوند: فصل ۵ مکانیزم‌های حفظ حریم خصوصی شامل شبه‌نام و تکنیک‌های بی‌نامی، فصل ۶ رمزنگاری و امضاهای دیجیتال، فصل ۷ مدیریت اعتماد و شهرت، فصل ۸ زنجیره بلوکی و \lr{IOTA Tangle} و فصل ۹ مدیریت هویت با شناسه‌های غیرمتمرکز و اعتبارنامه‌های تأییدپذیر را می‌پردازد. فصل ۱۰ نقش روش‌های هوش مصنوعی در تشخیص ناهنجاری و حملات را بررسی می‌کند. همچنین در فصل ۱۱ فناوری‌های نوظهور مانند شبکه‌های نسل پنجم و ششم، رایانش ابری، لبه و مه، و در فصل ۱۲ لایه سخت‌افزاری به عنوان بخش پشتیبان امنیت مورد توجه قرار می‌گیرند. در نهایت، فصل ۱۳ مسائل باز و جهت‌گیری‌های پژوهشی آینده و فصل ۱۴ جمع‌بندی کلی گزارش ارائه می‌شود.
